% Supplementary figures of the OPERA manuscript, shared with the thesis.

\clearpage
\setcounter{figure}{0}
\renewcommand{\thefigure}{S\arabic{figure}}
\setcounter{table}{0}
\renewcommand{\thetable}{S\arabic{table}}
\renewcommand{\resultfigureheight}{0.82\textheight}

\operasection{Supplementary Methods: model and analysis details}
% Detailed methods moved from the main Methods.
\subsection*{Sequence encoder and supervised training}
Our sequence models build on CORE-BEHRT \citep{odgaard2024core}, adapted to a
causal encoder with next-event pretraining and multi-outcome competing-risk
prediction. Each patient's history was represented as a chronological sequence
of diagnoses, procedures, medications, laboratory results (discretised into
value bins) and hospital contacts, preceded by sex and age. A causal transformer
encoder \citep{vaswani2017attention} with a hidden size of 256, four layers and
four attention heads, and a scaled Time2Vec encoding of time
\citep{kazemi2019time2vec}, was pretrained to predict the next event over a
context of 8,192 events. \placeholder{Number of individuals and events,
training budget and selected epoch.} This encoder initialised every supervised
model in the main comparisons.

\subsection*{Supervised training across diseases and outcomes}
Every supervised model consisted of the pretrained encoder, trained further,
and a linear head that produced, for each outcome, piecewise-constant hazards
for the outcome and for competing death, building on neural interval-hazard and
competing-risk models \citep{gensheimer2019nnetsurvival,lee2018deephit}.
Cumulative incidence at each horizon was derived from these hazards. The loss
was a binary cross-entropy at each horizon: patients with the event before the
horizon were positive, patients who died of another cause or were followed
beyond the horizon were negative, and patients censored earlier did not
contribute at that horizon. Losses were averaged within each outcome and then
across endpoints, where the grade variants of one laboratory finding (grade
$\geq$2 and $\geq$3) and nested definitions of the same finding formed one
endpoint (69 outcomes, 52 endpoints), so that a finding defined at two grades
did not count twice. This weighting is separate from the clinical domains used
to present results, which group related outcomes for display only.

The four supervision scopes differed only in which cells of the
disease-by-outcome grid---one outcome in one disease---provided labels
(\cref{fig:supervision-grid}); we use ``cell'' for such a combination
throughout. A specialist learned one outcome
in one disease; a disease-pooled model learned one outcome across all diseases;
an outcome-pooled model learned all outcomes within one disease; and the joint
model learned all outcomes in all diseases. Every model predicted all six
horizons and was trained for up to 30 epochs, stopping after five epochs
without improvement in tuning loss, with the same optimiser, batch size and
sampling settings, and every comparison used the same target patients, labels
and horizons. Because equal epochs do not mean equal optimisation steps or
computation, we report both for each scope.

\subsection*{Outcome panel}
We predicted 69 outcomes covering disease course (second-line treatment,
treatment failure, transformation, new malignancy and death), infection,
haematological toxicity and transfusion, kidney, liver and coagulation
toxicity, electrolyte and metabolic disturbances, cardiovascular events and
acute care, each at 30, 60, 90, 180, 365 and 730 days. Laboratory outcomes used
grading thresholds aligned with, but not identical to, the Common Terminology
Criteria for Adverse Events \citep{nci2025ctcae}, with baseline-aware thresholds
when pre-treatment values were already abnormal; infection was additionally
captured through antimicrobial treatment, which we label as such. A patient
contributed to an outcome only if it was ascertainable for them---for example,
if laboratory follow-up was available---and was otherwise treated as missing
rather than negative. Death was a competing event for every non-fatal outcome.
The 69 outcomes were selected from 87 candidates by a prespecified rule based on
training and tuning event counts at one year, without held-out data
(Supplementary Methods S8). Definitions, code lists and risk-set sizes are
given in Supplementary Methods S5--S9.

\subsection*{Tabular comparators: features and tuning}
For each disease--outcome cell and horizon, we trained gradient-boosted trees
(XGBoost; \citealp{chen2016xgboost}) and ridge-penalised logistic regression on the
same patients, labels and competing-event handling as the neural models. Their
features came from the pipeline of our earlier DLBCL study
\citep{werling2026dlbcl}: using timeseriesflattener
\citep{bernstorff2023tsflattener}, prescriptions, laboratory results,
procedures and referrals, pathology codes, diagnoses, blood cultures and
microbiology were summarised over 90-day, one-year and three-year windows before
the index date, with aggregates suited to each source (counts and distinct
counts for codes; latest, minimum, maximum, variance and slope for laboratory
values). Variables recorded for fewer than 10\% of training patients were
represented by per-source counts of such rare occurrences. To give the tabular
and sequence models the same information, quality-registry variables such as
stage and performance status were removed, and age and sex at the index date
were added. Unlike the earlier study, no LASSO preselection was applied.
Numeric features were median-imputed with missingness indicators and
standardised, and categorical features one-hot encoded, with all preprocessing
fitted on training data. Hyperparameters were chosen for each cell by log loss
on the tuning set from a prespecified grid: for XGBoost, shallow trees (depth
1--3), deeper trees (depth 4 and 6) and the setting of the earlier study
(depth 8, learning rate 0.01, up to 3,000 trees), with the number of trees also
chosen on the tuning set; for logistic regression, the penalty strength along
a seven-value path ($C$ from $3\times10^{-4}$ to 1, fitted with warm starts). Neither model was class-reweighted, so both produce
probabilities that can be compared directly. \placeholder{Confirm against the
delivered feature export: anchoring at first-line treatment, source list and
number of features.}

As a reference outside the information-matched comparison, XGBoost was also
trained for each cell with the clinical quality-registry variables recorded at
diagnosis added to the same features (performance status, stage, lactate
dehydrogenase, extranodal involvement, B symptoms and cytogenetic risk where the
registry for the disease records them; \placeholder{final list}). Values were
taken as recorded at diagnosis, left missing where not recorded, and handled by
the same missingness indicators; hyperparameters were tuned on the same grid.
The sequence models never saw these variables, so this arm estimates what the
curated registry adds, not a like-for-like comparison.

\subsection*{Comparison with clinical prognostic indices}
The quality registries record, for each patient, the prognostic index scored at
diagnosis (for example the IPI, FLIPI, MIPI, CLL-IPI or ISS) and its risk
category (low, intermediate or high), before treatment started. Every index
recorded for at least 50 patients of a subtype, and every index used in more
than one subtype pooled across all subtypes recording it, was compared with the joint
model's predicted two-year risk of treatment failure, with death part of the
endpoint, in the patients who had that index. For DLBCL we also computed the
NCCN-IPI from the components recorded in the lymphoma registry, following
the original definition \citep{zhou2014nccnipi}: age (41--60, 61--75 and
over 75 years: 1, 2 and 3 points), LDH relative to the upper limit of normal
(205~U/L, or 255~U/L from age 70; ratio above 1 and above 3: 1 and 2 points),
Ann Arbor stage III--IV (1 point), involvement of bone marrow, central nervous
system, liver, gastrointestinal tract or lung (1 point), and performance
status $\geq$2 (1 point); low 0--1, low-intermediate 2--3, high-intermediate
4--5, high $\geq$6. Patients missing any component were not scored.

Model risk groups were fixed on tuning patients with known two-year status and
applied unchanged to the test set. The low cutoff was the highest predicted
risk at which the group below it, of at least 20 patients, had an observed risk
no higher than the index's low-risk category. The high cutoff was the predicted
risk at which the model's specificity equalled that of the index's high-risk
category, so that both flagged the same proportion of patients without the
event; this is also the operating point marked in the precision--recall
comparison. Patients between the two cutoffs were divided in the proportions
of the index's intermediate categories (one for the registry indices, two for
the NCCN-IPI). Anchoring both ends to the index makes the groups directly
comparable while leaving the model free to differ: a model that ranks patients
better places more patients at low risk at the same risk, and captures more
events at the same specificity. For each index and each system we report group
size, observed risk, and the sensitivity and specificity of classifying a
patient in that group or higher, and, for the model, the observed risk in its
highest-risk 5\% and 10\% of patients, which the index cannot resolve within
its top category. As a sensitivity analysis, model groups were instead given
exactly the index's group sizes. Discrimination used the index category (or
NCCN-IPI points) as an ordinal score.

\subsection*{Routing formulation and selection}
To test whether a joint model benefits from conditioning its predictions on
disease and outcome beyond ordinary shared learning, we compared the joint
model with a routed model trained on the same panel. In the routed model, the
patient representation
$h_i$ is adjusted for each target outcome $k$ by a weighted combination of
small residual adapters, one per disease and one per outcome:
\begin{equation}
 \begin{aligned}
 h_{ik}&=h_i+\lambda\sum_s\pi_{iks}\,c_s\,A_s(h_i),\\
 A_s(h)&=U_s\,\mathrm{GELU}(V_s h),
 \qquad c_s=\frac{e_s}{e_s+\tau}.
 \end{aligned}
 \label{eq:router}
\end{equation}
where $A_s$ is a low-rank adapter for source $s$, the weights $\pi_{iks}$ are
computed from the patient, the patient's disease and the target outcome, and
$c_s$ shrinks the contribution of sources with few training events $e_s$. The
adapted representation enters the same linear head as in the joint model. Each adapter was
updated only by the supervision of its own disease or outcome. The approach
relates to mixtures of experts and parameter-efficient adapters
\citep{jacobs1991adaptive,ma2018mmoe,houlsby2019adapters,pfeiffer2021adapterfusion}.

The router configuration was chosen on the tuning set in a short prespecified
sequence: which adapter banks to include (none, disease, outcome or both), then
the adapter rank, then the form of the routing query. A simpler configuration
was preferred unless a more complex one improved tuning performance beyond a
prespecified margin, and a routed model with no adapters was required to
match the joint model before differences between adapter banks were interpreted.
\placeholder{Insert the selected configuration.} Because a target outcome
already has its own prediction head, an outcome adapter bank must add value
beyond that head; route weights describe how the model allocates corrections
and are not interpreted as evidence of relationships between diseases.



\begin{figure}[!htbp]
  \centering
  \resizebox{0.96\textwidth}{!}{\begin{tikzpicture}[opera,x=1cm,y=1cm]
  % One target prediction, represented as four clearly separated stages.
  \node[eyebrow,anchor=west] at (0,0) {SOURCE-ANCHORED ROUTER};
  \node[labelsmall,anchor=west] at (0,-.36)
    {One patient $i$ and one target outcome $k$};

  \node[eyebrow] at (1.35,-1.02) {INPUTS};
  \node[eyebrow] at (5.02,-1.02) {ROUTE};
  \node[eyebrow] at (9.28,-1.02) {SOURCE BANK};
  \node[eyebrow] at (14.10,-1.02) {PREDICTION};
  \draw[divider] (3.03,-.88)--(3.03,-5.08);
  \draw[divider] (7.02,-.88)--(7.02,-5.08);
  \draw[divider] (11.94,-.88)--(11.94,-5.08);

  % Input identities are not source representations: the encoder supplies z.
  \node[stageblue,minimum width=2.18cm] (patient) at (1.38,-2.08)
    {patient state $z_i$};
  \node[stageteal,minimum width=2.18cm] (disease) at (1.38,-3.08)
    {fine disease $d_i$};
  \node[stagewarm,minimum width=2.18cm] (target) at (1.38,-4.08)
    {target outcome $k$};

  \node[stageblue,minimum width=2.48cm,minimum height=1.04cm]
    (router) at (5.03,-3.08) {learned weights\\$\pi_{i,k,s}$};
  \draw[flowblue] (patient.east) to[out=0,in=160] (router.north west);
  \draw[flowteal] (disease.east)--(router.west);
  \draw[flowwarm] (target.east) to[out=0,in=200] (router.south west);
  \node[captionnote,text width=3.40cm] at (5.03,-4.36)
    {Patient, disease and target context};

  % The three rows are categories, not a claim that only three sources exist.
  \node[stageblue,minimum width=3.58cm] (shared) at (9.45,-1.97)
    {$z_i$ \quad shared path};
  \node[stageteal,minimum width=3.58cm] (dsrc) at (9.45,-3.08)
    {$\sum_d\pi_{i,k,d}\,a_d(z_i)$ \quad disease};
  \node[stagewarm,minimum width=3.58cm] (osrc) at (9.45,-4.19)
    {$\sum_o\pi_{i,k,o}\,a_o(z_i)$ \quad outcome};
  \draw[flowblue] (router.east) to[out=25,in=180] (shared.west);
  \draw[flowteal] (router.east)--(dsrc.west);
  \draw[flowwarm] (router.east) to[out=-25,in=180] (osrc.west);

  \node[stage,minimum width=2.33cm] (mix) at (14.07,-2.56)
    {weighted residual\\$z_{i,k}$};
  \node[outcome,minimum width=2.33cm] (hazards) at (14.07,-4.02)
    {interval hazards\\$\lambda_{i,k,t}$};
  \draw[flowblue] (shared.east) to[out=0,in=155] (mix.north west);
  \draw[flowteal] (dsrc.east)--(mix.west);
  \draw[flowwarm] (osrc.east) to[out=0,in=205] (mix.south west);
  \draw[flowwarm] (mix.south)--(hazards.north);

  \draw[divider] (0,-5.34)--(15.84,-5.34);
  \node[captionnote,anchor=west,text width=15.40cm,align=left]
    at (.20,-5.68)
    {The shared path remains unscaled; route weights govern adapter residuals. The bank includes eligible borrowed sources, whose adapter parameters are detached.};
\end{tikzpicture}
}
  \caption{\textbf{Source-anchored disease--outcome routing.} For each patient
  and target outcome, the router mixes a shared representation with residual
  adapters indexed by fine disease and outcome. Its weights depend on the
  patient, disease and target; sources with no training support are masked and
  low-event sources are shrunk. An adapter is updated by its own eligible
  supervision, while borrowed adapter parameters are detached. The mixed
  representation enters a shared linear head for interval hazards and
  cumulative incidence. This is a model schematic, not an empirical route-weight
  or performance plot.}
  \label{fig:source-router}
\end{figure}


\subsection*{Hierarchy for effect heterogeneity}
The primary evidence is the paired effect in each disease--outcome cell with
its patient-bootstrap interval. Many cells are small, so their raw estimates
are noisy and rankings of individual cells are unreliable. We therefore also
fit a deliberately simple Bayesian hierarchical model whose purpose is partial
pooling: each cell's effect is the panel mean plus a disease effect, an outcome
effect and a cell-specific deviation, so that a noisy cell is shrunk towards
what similar cells show, by an amount the data determine. The model answers
two questions: how large is the typical effect, and how much does it vary
between diseases and between outcomes? The cell estimates are not treated as
independent: their sampling covariance, taken from the shared patient-bootstrap
draws, enters the likelihood, because the same patients contribute to every
outcome within a disease.

Whether effects depend on the amount of training data available for the
target is a separate question; we do not assume a linear relationship.
A low-dimensional spline still constrains smoothness and complexity. The cell effect is allowed to vary smoothly with log training
events through a cubic spline with few knots and a shrinkage prior, to regularise departures from a flat relationship; the prior still affects
the fitted shape and is assessed in sensitivity checks. As a transparent check
that does not depend on smoothing, cells are also grouped into prespecified
bins of training events (fixed before held-out evaluation) and the mean effect
is estimated within each bin. Within each
question, the model fits the independent effects directly (for pooling:
disease pooling, outcome pooling and their interaction); conditional and
combined contrasts are exact sums of these, so all reported pooling contrasts
are mutually consistent. Priors
are weakly informative on each metric's scale and checked with prior
predictive simulation; convergence and posterior predictive checks are
reported. Posterior probabilities of benefit, harm and practical equivalence
use a prespecified tolerance. Raw paired estimates are always shown alongside.
\placeholder{Insert locked priors, tolerances, support-bin edges and sampling
diagnostics.}


\FloatBarrier

\clearpage
\operasection{Supplementary Methods: evaluation and reporting}
% Evaluation methods in reported form; empirical tables are filled from the locked analysis.
\subsection*{Risk sets and paired comparisons}
The unit of evaluation was one disease--outcome--horizon combination. Its risk
set was determined by eligibility and ascertainment of the outcome, not by the
models compared. For each split we report eligible patients, outcome events,
competing deaths and patients whose status was unknown at the horizon. Each
paired comparison used the patients with predictions from both models, so
different contrasts may rest on different risk sets, and their denominators are reported. Cells in which a model could not be fitted or a metric could not be
estimated remain in the coverage inventory with the reason
(\cref{fig:supp-coverage}) and are not treated as zero effects.

Superiority among evaluable cells was kept separate from coverage of
additional cells. A pooled model could predict cells in which no specialist
could be fitted. These cells have absolute performance estimates but no paired
specialist comparison, and were not imputed into pooled averages. Outcome
definitions and the reconstruction of treatment lines are in the appended
data-definition supplement.

\subsection*{Clinical prognostic indices}
Each index was evaluated in the population and for the endpoint closest to
those it was developed for, at the same prediction time. Original and evaluated
endpoints and horizons are listed side by side, and a shorter available
horizon is not presented as a reproduction of a longer development horizon.
All methods in an index comparison used the same patients, those with the index
recorded at diagnosis, and how often it was recorded is reported by disease and
calendar period. Index categories were used as recorded. Because the sequence
and tabular models read the wider history, a gain over an index is a gain of
the whole pipeline and does not separate the model from the information it
receives.

Index categories are not individual probabilities. They were used for ranking
and stratification. Brier scores and calibration of an index would require a
probability mapping fitted on development data, which we did not use.
\placeholder{Index, version, population and availability inventory from the
secure analysis.}

\subsection*{Risk groups, age and operating points}
Continuous model risks were converted to groups fixed at baseline with cutoffs
selected on tuning data, while index groups kept their original definitions.
Kaplan--Meier estimates describe endpoints that include death, and
Aalen--Johansen estimates non-fatal endpoints with death as a competing event.
Group sizes, events and 95\% bootstrap intervals accompany every estimate.
Patients without a recorded index were not included in index comparisons.

In the age figures, patient counts, mean predicted risk and observed risk are
shown as separate rows on the same age bands and groups. The predicted row is
the sequence model's mean probability, also within index categories. The index itself supplies no probability. Estimates are suppressed for cells with fewer
than ten patients, while counts remain visible.

Operating points used thresholds fixed on tuning data, so the specificity
matched on tuning data is not exactly matched in the test period. We report
the achieved sensitivity, specificity, positive and negative predictive values
and patient counts. Confusion counts use patients with known status at the
horizon, and the number excluded by censoring is reported. Censoring-weighted quantities are labelled as such.

For the panel as a whole, every model was also evaluated at six decision
points, each fixed per model, cell and horizon on that model's tuning
predictions. The first flagged as many tuning patients as had the event, so
that every model raised the same number of alerts. The second and third
maximised Youden's index and the MCC. The remaining three were the lowest
thresholds with a tuning specificity of at least 80\%, 90\% and 95\%. The test
period then shows how far the achieved specificity drifts from its tuning
target. For each decision point we report the threshold, the confusion matrix,
sensitivity, specificity, positive and negative predictive values, MCC and the
share of patients flagged, with 95\% intervals from the shared patient
bootstrap and paired differences between the joint model and XGBoost, the
joint model and the specialist, and the routed and joint models. Confusion
counts below five are suppressed in exported tables, as required by the secure
environment. Cells with fewer than five tuning events or non-events have no
decision point.

\subsection*{Censoring, calibration and robustness}
Censoring-weighted and known-status estimates answer different questions and
are reported separately, and the known-status analysis is a prespecified
sensitivity analysis. For the censoring weights we report the censoring model,
the censoring survival at each horizon and the distribution of weights. Long
horizons with little follow-up are flagged where estimates are unstable.

Calibration is reported as calibration-in-the-large, the calibration slope
where estimable, and reliability curves for outcomes chosen by a fixed rule. No
model was recalibrated on test data. Sensitivity analyses by horizon, calendar
year, amount of recorded history and timing of the subtype diagnosis are
reported on their own risk sets. Raw paired effects are reported with the
hierarchical summaries, and the support curve, unsmoothed and binned analyses
are treated as views of the same data rather than independent replications.

\subsection*{Patient representations}
The representation figures show the same test patients at five stages: local
pretraining, the locally pretrained joint model, national pretraining, national
pretraining with local adaptation, and the nationally initialised joint model.
For each stage, the projection was fitted on tuning patients and applied
unchanged to test patients, so coordinates are comparable within a stage but
not between stages. For every stage we record the checkpoint, the layer and
pooling used, the input cutoff and truncation. Clusters and changes in
separation are descriptive. They do not show differences in prediction,
biological mechanisms or the effect of transfer. Only kernel-smoothed
summaries over at least \placeholder{N} patients are shown, following the
disclosure rules of the secure environment.

\subsection*{Numerical tables}
The supplementary tables give, for every cell, the risk sets, absolute
performance, paired effects, index comparisons, calibration, operating points
and sensitivity analyses, together with training resources. Each model was
trained once, so intervals describe sampling of test patients only. Readable
summaries accompany the complete cell-level files.
\placeholder{Populate from the locked analysis.}

\IfFileExists{tables/generated/supplement_tables.tex}{%
  \subsection*{Reviewed numerical summaries}
  \input{tables/generated/supplement_tables}
}{%
  \noindent\placeholder{Reviewed aggregate summary tables have not yet been
  exported. Complete cell-level numerical tables will accompany the figures.}
}

\clearpage
\operasection{Supplementary figures}

Supplementary figures group study design, performance, clinical comparisons,
representations and explanation diagnostics; all are cross-referenced in the text.

\begin{figure}[p]
  \centering
  \resizebox{\textwidth}{!}{% Supervision scopes drawn as four copies of the Fig. 2B disease-outcome matrix:
% same axes, rows/columns, ellipses, cell geometry, lavender gradient, sienna
% target and teal borrowing spokes. Filled cells are supervised tasks; pale
% outlined cells are present in the panel but not supervised by that fit.
% Lines are illustrative sharing paths, not learned gates or measured routes.
\begin{tikzpicture}[opera,x=1cm,y=1cm]
\foreach \scopeid/\sx/\sy/\letter/\eyebrowtext/\title/\sub/\note in {
 S/0/0/A/{SPECIALIST}/{Learns from its own cell}/{One disease, one outcome: one fit per supported cell.}/{The target learns only from its own labels.},
 D/10.75/0/B/{DISEASE-POOLED}/{Learns across diseases}/{All applicable diseases, one outcome: one fit per outcome.}/{The target borrows from the same outcome in other diseases.},
 O/0/-9.55/C/{OUTCOME-POOLED}/{Learns across outcomes}/{One disease, all applicable outcomes: one fit per disease.}/{The target borrows from related outcomes in the same disease.},
 J/10.75/-9.55/D/{JOINT}/{Learns from the whole panel}/{All diseases and outcomes: one shared fit. OPERA uses this same supervision with task-specific corrections.}/{Every supervised cell can contribute through the shared encoder and readout.}}{
\begin{scope}[xshift=\sx cm,yshift=\sy cm]
 \node[panelletter,anchor=north west] at (0,0) {\letter};
 \node[eyebrow,anchor=north west] at (.48,-.03) {\eyebrowtext};
 \node[heading,anchor=north west] at (.48,-.38) {\title};
 \node[labelsmall,anchor=north west,align=left,text width=8.6cm] at (.48,-.82) {\sub};
 \begin{scope}[yshift=-5mm]
 % Axes exactly as in Fig. 2B.
 \draw[-{Latex[length=1.8mm]},draw=OperaIndigo,line width=.75pt] (2.76,-1.55)--(9.12,-1.55);
 \node[eyebrow,text=OperaIndigo] at (5.94,-1.30) {MORE OUTCOME TASKS};
 \foreach \r/\disease in {0/DLBCL,1/Multiple myeloma,2/CLL,3/Hodgkin lymphoma,4/PTCL}
  \node[labelsmall,anchor=east] at (2.60,-2.72-.84*\r) {\disease};
 \draw[-{Latex[length=1.8mm]},draw=OperaTeal,line width=.75pt] (.16,-2.40)--(.16,-7.62);
 \node[eyebrow,text=OperaTeal,rotate=90] at (-.13,-5.01) {MORE DISEASE COHORTS};
 \foreach \col/\outcome in {0/Anaemia,1/Infection,2/Renal,3/Fracture,4/Survival}
  \node[font=\sffamily\fontsize{6.2}{7}\selectfont,text=OperaInkSecondary,
   align=center,text width=1.05cm,anchor=south] at (3.22+1.12*\col,-2.27) {\outcome};
 \node[font=\sffamily\fontsize{6.2}{7}\selectfont,text=OperaInkSecondary,
  align=center,text width=1.05cm,anchor=south] at (9.12,-2.27) {Outcome$_n$};
 % Main 5 x 5 block. Supervision depends on the scope; the target is row 4, column 1.
 \foreach \r in {0,...,4}{\foreach \col in {0,...,4}{
  \def\active{0}
  \if\scopeid D \ifnum\col=1 \def\active{1}\fi\fi
  \if\scopeid O \ifnum\r=4 \def\active{1}\fi\fi
  \if\scopeid J \def\active{1}\fi
  \def\tone{OperaPaper}
  \ifnum\active=1 \pgfmathtruncatemacro{\shade}{7+3*\r+2*\col}\edef\tone{OperaIndigo!\shade}\fi
  \node[rounded corners=.35mm,draw=OperaBorder,fill=\tone,
   minimum width=8.8mm,minimum height=5.7mm] (\scopeid m\r\col) at (3.22+1.12*\col,-2.72-.84*\r) {};
 }}
 \foreach \r in {0,...,4}{
  \node[font=\sffamily\scriptsize,text=OperaInkMuted] at (8.41,-2.72-.84*\r) {$\cdots$};
  \def\active{0}
  \if\scopeid O \ifnum\r=4 \def\active{1}\fi\fi
  \if\scopeid J \def\active{1}\fi
  \def\tone{OperaPaper}
  \ifnum\active=1 \pgfmathtruncatemacro{\shade}{16+3*\r}\edef\tone{OperaIndigo!\shade}\fi
  \node[rounded corners=.35mm,draw=OperaBorder,fill=\tone,
   minimum width=8.8mm,minimum height=5.7mm] (\scopeid n\r) at (9.12,-2.72-.84*\r) {};
 }
 \foreach \x in {3.22,4.34,5.46,6.58,7.70,9.12}{
  \node[font=\sffamily\scriptsize,text=OperaInkMuted] at (\x,-6.60) {$\vdots$};
 }
 \node[font=\sffamily\scriptsize,text=OperaInkMuted] at (8.41,-6.60) {$\ddots$};
 \node[labelsmall,anchor=east] at (2.60,-7.29) {Disease$_m$};
 \foreach \col in {0,...,4}{
  \def\active{0}
  \if\scopeid D \ifnum\col=1 \def\active{1}\fi\fi
  \if\scopeid J \def\active{1}\fi
  \def\tone{OperaPaper}
  \ifnum\active=1 \pgfmathtruncatemacro{\shade}{22+2*\col}\edef\tone{OperaIndigo!\shade}\fi
  \node[rounded corners=.35mm,draw=OperaBorder,fill=\tone,
   minimum width=8.8mm,minimum height=5.7mm] (\scopeid z\col) at (3.22+1.12*\col,-7.29) {};
 }
 \def\tone{OperaPaper}
 \if\scopeid J \def\tone{OperaIndigo!34}\fi
 \node[rounded corners=.35mm,draw=OperaBorder,fill=\tone,
  minimum width=8.8mm,minimum height=5.7mm] (\scopeid zn) at (9.12,-7.29) {};
 \node[rounded corners=.35mm,draw=OperaSienna,line width=1.1pt,fill=OperaSienna!8,
  minimum width=8.8mm,minimum height=5.7mm] (\scopeid target) at (4.34,-6.08) {};
 % Spokes from supervised cells only: bold when the cell shares the target's
 % outcome (column 1) or disease (row 4), thin otherwise, as in Fig. 2B.
 \begin{scope}[on background layer]
  \foreach \r in {0,...,4}{\foreach \col in {0,...,4}{
   \def\link{0}
   \if\scopeid D \ifnum\col=1 \def\link{1}\fi\fi
   \if\scopeid O \ifnum\r=4 \def\link{1}\fi\fi
   \if\scopeid J \def\link{1}\fi
   \ifnum\r=4 \ifnum\col=1 \def\link{0}\fi\fi
   \ifnum\link=1
    \def\bold{0}
    \ifnum\col=1 \def\bold{1}\fi
    \ifnum\r=4 \def\bold{1}\fi
    \ifnum\bold=1
     \draw[draw=OperaTeal,opacity=1,line width=1.6pt] (\scopeid m\r\col)--(\scopeid target);
    \else
     \draw[draw=OperaTeal,opacity=.55,line width=.6pt] (\scopeid m\r\col)--(\scopeid target);
    \fi
   \fi
  }}
  \if\scopeid O \draw[draw=OperaTeal,opacity=1,line width=1.6pt] (\scopeid n4)--(\scopeid target);\fi
  \if\scopeid D \draw[draw=OperaTeal,opacity=1,line width=1.6pt] (\scopeid z1)--(\scopeid target);\fi
  \if\scopeid J
   \foreach \r in {0,...,3}{\draw[draw=OperaTeal,opacity=.55,line width=.6pt] (\scopeid n\r)--(\scopeid target);}
   \draw[draw=OperaTeal,opacity=1,line width=1.6pt] (\scopeid n4)--(\scopeid target);
   \foreach \col in {0,2,3,4}{\draw[draw=OperaTeal,opacity=.55,line width=.6pt] (\scopeid z\col)--(\scopeid target);}
   \draw[draw=OperaTeal,opacity=1,line width=1.6pt] (\scopeid z1)--(\scopeid target);
   \draw[draw=OperaTeal,opacity=.55,line width=.6pt] (\scopeid zn)--(\scopeid target);
  \fi
 \end{scope}
 \node[captionnote,text width=8.5cm] at (5.60,-7.98) {\note};
 \end{scope}
\end{scope}
}
\draw[divider] (10.05,.05)--(10.05,-18.85);
\draw[divider] (-.35,-9.15)--(20.15,-9.15);
\node[labelsmall,text=OperaTeal,align=center,text width=18cm] at (9.90,-19.30)
 {Sienna outline: the same target task in every panel. Filled cells: supervised tasks; pale cells: not supervised by that fit.
 Thin and bold lines are illustrative sharing paths, bold where the cell shares the target's disease or outcome.};
\end{tikzpicture}
}
  \caption{\textbf{Supervision scopes on the disease--outcome matrix.} The matrix of \cref{fig:paradigms}E, with one target task
  (outlined) evaluated identically in every panel. Diseases are rows and
  outcomes are columns. \textbf{A}, A specialist learns one outcome in one disease. \textbf{B}, A disease-pooled model learns one outcome across all diseases. \textbf{C}, An outcome-pooled model learns all outcomes within one disease. \textbf{D}, The joint model learns all outcomes in all diseases.
  Filled cells indicate supervision, not performance. The routed model uses
  the same panel as the joint model and changes only how predictions are read
  out. Every model predicts all six horizons.}
  \label{fig:supervision-grid}
\end{figure}

\begin{figure}[p]
\centering
\resizebox{\textwidth}{!}{% What the sequence model reads: one card per source (A) and how a patient's
% history becomes one token sequence cut at the start of treatment (B).
% Requires \usepackage{fontawesome5} and opera_style.tex.
\begin{tikzpicture}[opera,x=1cm,y=1cm,
    src/.style={rounded corners=0.8mm, draw=#1!45, fill=#1!6, line width=.5pt,
      minimum width=10.1cm, minimum height=1.36cm},
    icon/.style={circle, fill=#1, text=white, minimum size=7.6mm, inner sep=0pt,
      font=\normalsize},
    srctitle/.style={font=\sffamily\fontsize{8}{9}\selectfont\bfseries, text=OperaInk,
      anchor=north west},
    srcbody/.style={font=\sffamily\fontsize{6.6}{7.8}\selectfont, text=OperaInkSecondary,
      anchor=north west, align=left, text width=5.35cm},
    srctok/.style={token, draw=#1!55, fill=white, font=\ttfamily\fontsize{5.4}{6}\selectfont,
      minimum height=4.2mm, inner xsep=1mm},
    srctime/.style={font=\sffamily\fontsize{6}{7}\selectfont, text=OperaInkMuted,
      anchor=north east, align=right},
    seqtok/.style={rounded corners=.5mm, fill=#1, minimum width=3.3mm, minimum height=5.2mm,
      inner sep=0pt},
  ]
  % ---------------------------------------------------------------- A
  \node[panelletter,anchor=north west] at (0,0) {A};
  \node[eyebrow,anchor=north west] at (.45,-.03) {DATA SOURCES};
  \node[heading,anchor=north west] at (.45,-.36) {Registers and example tokens};

  % One card per source: y, colour, icon, title, contents and timing, two example tokens.
  \def\sourcecard#1#2#3#4#5#6#7{%
    \node[src=#2] at (5.05,#1) {};
    \node[icon=#2] at (.62,#1) {\faIcon{#3}};
    \node[srctitle] at (1.2,#1+.55) {#4};
    \node[srcbody] at (1.2,#1+.23) {#5};
    \node[srctok=#2,anchor=west] at (6.75,#1+.22) {#6};
    \node[srctok=#2,anchor=west] at (6.75,#1-.25) {#7};}
  \sourcecard{-1.55}{OperaSlate}{user}{Demographics}
    {Sex, and age at the start of treatment.}
    {SEX//F}{AGE//71}
  \sourcecard{-3.0}{OperaIndigo}{hospital}{National patient register}
    {Every hospital contact: diagnoses (ICD-10) at discharge; procedures and the
     contact's type, reason, priority, specialty and region at admission.}
    {DIAGNOSIS//DC833}{CONTACT\_SPECIALTY//HAEM}
  \sourcecard{-4.45}{OperaLavender}{microscope}{Pathology register}
    {Diagnoses from tissue and cell samples (SNOMED) and the requesting
     department, at requisition.}
    {PATHOLOGY//M96803}{REGION//HOVEDSTADEN}
  \sourcecard{-5.9}{OperaTeal}{prescription-bottle-alt}{Prescription register}
    {Every prescription redeemed at a community pharmacy (ATC), at
     dispensing.}
    {MEDICATION//J01CA04}{MEDICATION//N02BE01}
  \sourcecard{-7.35}{OperaOchre}{syringe}{Hospital medication register}
    {Drugs given in hospital, including chemotherapy (ATC), and the treating
     department, at the start of the order.}
    {MEDICATION//L01XC02}{CONTACT\_SPECIALTY//ONC}
  \sourcecard{-8.8}{OperaSienna}{vial}{Laboratory database}
    {Laboratory results (NPU codes) at sampling: the test with its value bin,
     or the text result.}
    {LABTEST//NPU02319//BIN\_7}{LABTEST//NPU19748//NEG}
  \draw[divider] (10.55,.05)--(10.55,-9.5);

  % ---------------------------------------------------------------- B
  \node[panelletter,anchor=north west] at (10.9,0) {B};
  \node[eyebrow,anchor=north west] at (11.35,-.03) {INPUT AND OUTPUT};
  \node[heading,anchor=north west] at (11.35,-.36) {Patient sequence and predictions};

  % Timeline with tokens coloured by source; the cut at treatment start.
  \draw[-{Latex[length=1.6mm]},draw=OperaInkMuted,line width=.55pt] (11.1,-2.25)--(16.75,-2.25);
  \node[labelsmall,anchor=north east] at (16.75,-2.32) {time};
  \foreach \x/\col in {11.25/OperaSlate,11.62/OperaSlate}
    \node[seqtok=\col] at (\x,-1.85) {};
  \foreach \x/\col in {12.15/OperaIndigo,12.5/OperaTeal,12.85/OperaSienna,
      13.2/OperaSienna,13.55/OperaIndigo,13.9/OperaTeal,14.25/OperaLavender,
      14.6/OperaIndigo,14.95/OperaSienna,15.3/OperaSienna,15.65/OperaOchre,
      16.0/OperaSienna}
    \node[seqtok=\col] at (\x,-1.85) {};
  \draw[OperaSienna,line width=1.1pt] (16.32,-1.35)--(16.32,-2.55);
  \node[font=\sffamily\fontsize{6.4}{7}\selectfont\bfseries,text=OperaSienna,anchor=south]
    at (16.32,-1.33) {treatment starts};
  \node[labelsmall,anchor=south] at (11.44,-1.55) {sex, age};
  \draw[decorate,decoration={brace,amplitude=3pt,mirror},draw=OperaInkMuted]
    (12.0,-2.62)--(16.05,-2.62);
  \node[labelsmall,anchor=north,text width=4.3cm] at (14.08,-2.78)
    {everything recorded before treatment, in time order (up to 8,192 tokens)};

  \node[stageblue,minimum width=5.2cm] (enc) at (13.95,-4.4)
    {\textbf{Causal transformer encoder}\\[-1pt]
     {\fontsize{6.6}{7.6}\selectfont pretrained to predict the next event}};
  \draw[flowblue] (13.95,-3.45)--(enc.north);
  \node[softpanel,draw=OperaIndigo!42,fill=OperaIndigo!4,
    minimum width=3.6cm,minimum height=6mm,font=\sffamily\fontsize{6.6}{7.6}\selectfont]
    (z) at (13.95,-5.6) {shared representation $z$};
  \draw[flowblue,shorten >=1pt] (enc.south)--(z.north);

  % Outcomes as in Figure 1C: one probability per outcome and horizon.
  \node[labelsmall] (c1) at (11.75,-6.75) {$P($failure by $t)$};
  \node[labelsmall] (c2) at (13.5,-6.75) {$P($infection by $t)$};
  \node[font=\sffamily\scriptsize,text=OperaInkMuted] at (14.78,-6.75) {$\cdots$};
  \node[labelsmall] (cn) at (16.05,-6.75) {$P($outcome$_{69}$ by $t)$};
  \foreach \c in {c1,c2,cn}
    \draw[flowblue,shorten >=1pt] (z)--(\c.north);
  \begin{scope}[shift={(11.27,0)}]
    \foreach \i in {0,...,11}{
      \pgfmathsetmacro{\xa}{0.08*\i}
      \pgfmathsetmacro{\xb}{0.08*(\i+1)}
      \pgfmathtruncatemacro{\tone}{6+7*\i}
      \fill[OperaIndigo!\tone] (\xa,-7.27) rectangle (\xb,-6.99);
    }
  \end{scope}
  \begin{scope}[shift={(13.02,0)}]
    \foreach \i in {0,...,11}{
      \pgfmathsetmacro{\xa}{0.08*\i}
      \pgfmathsetmacro{\xb}{0.08*(\i+1)}
      \pgfmathtruncatemacro{\tone}{6+7*\i}
      \fill[OperaIndigo!\tone] (\xa,-7.27) rectangle (\xb,-6.99);
    }
  \end{scope}
  \begin{scope}[shift={(15.57,0)}]
    \foreach \i in {0,...,11}{
      \pgfmathsetmacro{\xa}{0.08*\i}
      \pgfmathsetmacro{\xb}{0.08*(\i+1)}
      \pgfmathtruncatemacro{\tone}{6+7*\i}
      \fill[OperaIndigo!\tone] (\xa,-7.27) rectangle (\xb,-6.99);
    }
  \end{scope}
  \node[labelsmall,anchor=north] at (13.95,-7.55)
    {$t$ = 30, 60, 90, 180, 365 and 730 days};
  \node[captionnote,anchor=north,text width=5.4cm] at (13.95,-8.35)
    {The same encoder serves every disease and outcome; only the labels each model is
     trained on differ.};
\end{tikzpicture}
}
\caption{\textbf{Data sources and input representation of the sequence model.} \textbf{A}, the six sources linked
through the civil registration number, what each contributes and when its events
are dated, with two example tokens per source. Codes are illustrative, and laboratory values enter as quantile bins fitted on events before 2022. \textbf{B}, each patient's history becomes one token sequence, coloured by source,
in time order and ending at the start of first-line treatment. The pretrained
encoder reads it once into a shared representation, from which the model gives
the probability of every outcome by each of six horizons.}
\label{fig:model-inputs}
\end{figure}

\begin{figure}[p]
  \centering
  \resizebox{0.94\textwidth}{!}{\begin{tikzpicture}[opera,node distance=10mm and 12mm,
    every node/.append style={execute at begin node={\hyphenpenalty=10000\exhyphenpenalty=10000}}]
  \node[stageblue,text width=6.1cm] (pretrainbox)
    {78,668 diagnosis-linked patients\\{\scriptsize register histories before 2022 used for encoder pretraining}};
  \node[eyebrow,above=3mm of pretrainbox] (pretrainlabel) {PRETRAINING POPULATION};
  \node[eyebrow,below=7mm of pretrainbox] (branchlabel)
    {FIRST-LINE COHORT ASCERTAINMENT};
  \node[stage,text width=6.1cm,below=4mm of branchlabel] (index)
    {First-line treatment index identified\\38,635 patients};
  \node[stage,text width=6.1cm,below=of index] (pass1)
    {Initial disease assignment using diagnoses through day 14\\37,357 assigned; 1,278 unresolved};
  \node[stageteal,text width=6.1cm,below=of pass1] (pass2)
    {Coding-lag reassessment through day 90\\209 additional patients assigned};
  \node[stage,text width=6.1cm,below=of pass2] (classified)
    {Disease assigned after both passes\\37,566 patients};
  \node[stageblue,text width=6.1cm,below=of classified] (firstline)
    {First-line cohort after IgM MGUS exclusion\\37,557 patients};
  \node[stageblue,text width=6.1cm,below=of firstline] (supervised)
    {37,531 patients for training, tuning and testing};

  \foreach \a/\b in {index/pass1,pass1/pass2,
                       pass2/classified,classified/firstline,firstline/supervised}
    \draw[flowblue,shorten >=1pt] (\a.south) -- (\b.north);
  \draw[flowblue,shorten >=1pt] (pretrainbox.south) -- (index.north);
  \node[eyebrow,fill=white,inner xsep=2mm] at (branchlabel) {FIRST-LINE COHORT ASCERTAINMENT};
  \node[eyebrow,fill=white,inner xsep=2mm,above=2mm of supervised]
    {FINAL SUPERVISED COHORT};

  % Mutually distinct exits from the central first-line flow.
  \node[softpanel,text width=4.45cm,align=left,right=of pretrainbox] (noindex)
    {\bfseries Excluded: \textit{n}\,=\,40{,}033\\No qualifying first-line treatment index};
  \node[softpanel,text width=4.45cm,align=left,right=of pass2] (unresolved)
    {\bfseries Excluded: \textit{n}\,=\,1{,}069\\Disease unresolved after both passes};
  \node[softpanel,text width=4.45cm,align=left,right=of classified] (secondary)
    {\bfseries Excluded: \textit{n}\,=\,9\\IgM MGUS, untreated by definition};
  \node[softpanel,text width=4.45cm,align=left,right=of firstline] (nosurvival)
    {\bfseries Excluded: \textit{n}\,=\,26\\Died (\textit{n}\,=\,17) or follow-up ended (\textit{n}\,=\,9) before treatment};
  \draw[relation,draw=OperaSienna,shorten >=1pt] (pretrainbox.east) -- (noindex.west);
  \draw[relation,draw=OperaSienna,shorten >=1pt] (pass2.east) -- (unresolved.west);
  \draw[relation,draw=OperaSienna,shorten >=1pt] (classified.east) -- (secondary.west);
  \draw[relation,draw=OperaSienna,shorten >=1pt] (firstline.east) -- (nosurvival.west);

  \node[stage,text width=3.55cm,below left=14mm and 14mm of supervised] (train)
    {Training\\index through 2021\\\textit{n}\,=\,29{,}105};
  \node[stage,text width=3.55cm,below=14mm of supervised] (tune)
    {Tuning\\index in 2022\\\textit{n}\,=\,1{,}980};
  \node[stage,text width=3.55cm,below right=14mm and 14mm of supervised] (held)
    {Test (held out)\\index from 2023\\\textit{n}\,=\,6{,}446};
  \coordinate (split) at ($(supervised.south)+(0,-7mm)$);
  \draw[flowblue] (supervised.south) -- (split);
  \draw[draw=OperaIndigo,line width=0.75pt]
    (train.north |- split) -- (held.north |- split);
  \foreach \dest in {train,tune,held}
    \draw[flowblue,shorten >=1pt] (\dest.north |- split) -- (\dest.north);

\end{tikzpicture}
}
  \caption{\textbf{CONSORT diagram.} Of 78,668 adults in DALY-CARE with a
  lymphoid cancer or related diagnosis, whose pre-2022 histories were used for
  pretraining, 38,635 started first-line treatment, 37,557 had an assigned lymphoid
  disease, and 37,531 alive and followed at treatment start were assigned to the training (to 2021),
  tuning (2022) or test (from 2023) period. Subtypes
  were assigned from diagnoses up to 14 days after treatment start, with
  reassessment up to day 90 for late-coded diagnoses. Counts are unique patients. Not every patient contributes to every endpoint: outcome-specific eligibility,
  competing events and follow-up are applied within each split, and the
  resulting risk sets are given in the Supplementary Methods.}
  \label{fig:cohort-flow}
\end{figure}

\begin{figure}[p]
  \resultfigure[0.9\textwidth]{supp_coverage}{Estimability}{Estimability of each
    disease--outcome cell.}
  \caption{\textbf{Estimability of disease--outcome cells.} Status of every
  disease--outcome cell at one year: all models evaluated, too few training
  events to train a specialist, no events in the test period, fit failed, or
  not applicable (sex- or disease-specific outcomes). Counts in the legend.}
  \label{fig:supp-coverage}
\end{figure}

\begin{figure}[p]
  \resultfigure[0.9\textwidth]{supp_support}{Training support}{Training
    events per cell and their cumulative distribution.}
  \caption{\textbf{Training events per disease--outcome cell.}
  \textbf{A}, Training events by one year in each disease--outcome cell (log scale, grey for none). \textbf{B}, Cumulative distribution across cells, with
  the support-bin edges used in the analysis marked.}
  \label{fig:supp-support}
\end{figure}

\begin{figure}[p]
  \resultfigure[0.95\textwidth]{supp_atlas_matrix}{Every outcome in every disease}
    {Test-period AUC of the joint model for each outcome in each disease.}
  \caption{\textbf{Discrimination of the joint model for each outcome in each
  disease.} \textbf{A}, Test-period
  one-year AUC of the joint model for each outcome (rows, by clinical domain)
  in each disease (columns, by disease group). Grey cells had too few events to estimate. \textbf{B}, The same outcomes over all diseases at each horizon.}
  \label{fig:supp-atlas-matrix}
\end{figure}

\begin{figure}[p]
  \resultfigure[0.9\textwidth]{supp_pooling_cells}{Cell-level pooling
    gains}{Disease by outcome maps of each pooling gain.}
  \caption{\textbf{Pooling gains for each disease--outcome cell.} As
  \cref{fig:pooling}A--D at the resolution of single cells. Diseases (rows) run from fewest to most training patients and outcomes (columns) from fewest to most training events, with log-scale bars showing both. Dots mark cells whose 95\% interval excludes zero: white for a gain, an open ring for a loss. Most cells are small, so colours alone do not establish benefit or harm. Dots mark unadjusted
  pointwise intervals, not multiplicity-corrected discoveries. Grey cells could not be estimated.}
  \label{fig:supp-pooling-cells}
\end{figure}

\begin{figure}[p]
  \resultfigure[0.9\textwidth]{supp_horizons}{Pooling gains across
    horizons}{Gains at 30 to 730 days.}
  \caption{\textbf{Pooling gains across horizons.} Median paired AUC gain
  over the specialist by clinical domain at each horizon for the
  disease-pooled, outcome-pooled and joint models.}
  \label{fig:supp-horizons}
\end{figure}

\begin{figure}[p]
  \resultfigure{fig_calibration}{Calibration}
    {A: calibration slope by clinical domain. B: calibration in the large.
    C: reliability curves for four example outcomes.}
  \caption{\textbf{Calibration of the joint, routed, specialist and XGBoost
  models.} \textbf{A},
  Calibration slope at one year for the joint, routed, specialist and
  XGBoost models by clinical domain (median and interquartile range across cells, ideal 1). \textbf{B}, Calibration in the large, observed minus
  mean predicted risk (ideal 0). \textbf{C}, Reliability curves for four example outcomes. Observed risks are weighted by the inverse probability of censoring and point size reflects the number of patients per bin. No model
  was recalibrated on test data.}
  \label{fig:calibration}
\end{figure}

\begin{figure}[p]
  \resultfigure{supp_regions}{Performance by region}
    {Gain of the joint model by the region where patients were treated.}
  \caption{\textbf{Performance by region.} The models were trained on patients from all five regions, and test patients are grouped by the
  region of most of their pre-treatment contacts and every comparison is
  repeated within each region on the same fits. Maps show \textbf{A}, the one-year AUC gain of the joint model over per-task XGBoost, \textbf{B}, the joint model's one-year AUC, and \textbf{C}, its calibration in the large (observed minus predicted risk). Below, the same with 95\% paired bootstrap intervals over
  equally weighted cells: \textbf{D}, gains over the specialist and XGBoost, \textbf{E}, AUC of the joint model and XGBoost, \textbf{F}, calibration in the large, and \textbf{G}, recorded events and laboratory results before
  treatment (median and interquartile range), which may explain regional
  differences. Region boundaries from geoBoundaries (CC BY 4.0)
  \citep{runfola2020geoboundaries}, with Bornholm as an inset.}
  \label{fig:supp-regions}
\end{figure}

\begin{figure}[p]
  \resultfigure{fig_clinical_score}{Operating points relative to the IPI}
    {A: precision--recall with matched-specificity operating point. B: model
    risk within index groups. C--D: incidence by age and risk group.}
  \caption{\textbf{Operating points and reclassification relative to the IPI
  in DLBCL.} \textbf{A},
  Precision--recall curves for death within two years (patients with known status at two years) for the model and for each IPI category
  threshold. The square marks the IPI high-risk category. Dotted lines project it onto the model's curve: straight up gives the model's precision at the IPI's recall, and straight right its recall at the IPI's precision. The diamond marks the model at a threshold chosen on tuning data to match the IPI's specificity. AP, average
  precision. \textbf{B}, Predicted two-year risk of each patient within each
  IPI category, coloured by whether the patient died, with the model risk groups shaded. \textbf{C--D}, Two-year overall survival (95\% bootstrap
  interval) and number of patients by age band and risk group for the model
  (\textbf{C}) and the IPI (\textbf{D}). Counts below five are shown as $<$5,
  and estimates are suppressed for cells with fewer than ten patients.}
  \label{fig:clinical-score}
\end{figure}

\begin{figure}[p]
  \resultfigure{supp_clinical_score_age}{Risk groups by age}{Patient counts, predictions and observed risk separately.}
  \caption{\textbf{Model and IPI risk groups by age in DLBCL.} Patients with DLBCL and a
  recorded IPI, grouped by model risk group (left) and IPI category (right).
  Rows show, by age band, the number of patients, the mean predicted two-year overall survival and the observed Kaplan--Meier estimate, each
  with its 95\% bootstrap interval. Estimates are suppressed for cells with fewer than ten patients, but counts remain visible.}
  \label{fig:clinical-score-age}
\end{figure}

\begin{figure}[p]
  \resultfigure[0.95\textwidth]{supp_score_overview}{Comparison with each recorded prognostic index}
    {Survival by model risk group and by index category for each index and
    population.}
  \caption{\textbf{Model risk groups compared with each recorded prognostic
  index.} One row
  per index and the population it applies to, on patients with that index
  recorded at diagnosis. Left, survival (or freedom from the endpoint) by model risk group. Middle, the same by index category, on the same patients. Right, estimates at two years with 95\% CI (M, model, and I, index). Shaded bands are 95\%
  bootstrap intervals and the dashed line marks two years. Model cutoffs were
  fixed on tuning data for each index, matching its low-risk observed risk and
  high-risk specificity. Every index recorded for at least 50 patients of a
  subtype is shown, and the NCCN-IPI in DLBCL. Group sizes, observed risk,
  sensitivity and specificity of both systems are in Supplementary
  Table~\placeholder{S}. \placeholder{Endpoint per index.}}
  \label{fig:score-overview}
\end{figure}

% One plate like the figure above for every recorded index, written by the
% clinical-scores step; copy its clinical_scores/ folder into the figure directory.
\IfFileExists{\resultfiguredir clinical_scores/supplement_clinical_scores.tex}{%
  \input{\resultfiguredir clinical_scores/supplement_clinical_scores.tex}}{}

\begin{figure}[p]
  \centering
  \resizebox{0.96\textwidth}{!}{\begin{tikzpicture}[opera,x=1cm,y=1cm]
  % One target prediction, represented as four clearly separated stages.
  \node[eyebrow,anchor=west] at (0,0) {SOURCE-ANCHORED ROUTER};
  \node[labelsmall,anchor=west] at (0,-.36)
    {One patient $i$ and one target outcome $k$};

  \node[eyebrow] at (1.35,-1.02) {INPUTS};
  \node[eyebrow] at (5.02,-1.02) {ROUTE};
  \node[eyebrow] at (9.28,-1.02) {SOURCE BANK};
  \node[eyebrow] at (14.10,-1.02) {PREDICTION};
  \draw[divider] (3.03,-.88)--(3.03,-5.08);
  \draw[divider] (7.02,-.88)--(7.02,-5.08);
  \draw[divider] (11.94,-.88)--(11.94,-5.08);

  % Input identities are not source representations: the encoder supplies z.
  \node[stageblue,minimum width=2.18cm] (patient) at (1.38,-2.08)
    {patient state $z_i$};
  \node[stageteal,minimum width=2.18cm] (disease) at (1.38,-3.08)
    {fine disease $d_i$};
  \node[stagewarm,minimum width=2.18cm] (target) at (1.38,-4.08)
    {target outcome $k$};

  \node[stageblue,minimum width=2.48cm,minimum height=1.04cm]
    (router) at (5.03,-3.08) {learned weights\\$\pi_{i,k,s}$};
  \draw[flowblue] (patient.east) to[out=0,in=160] (router.north west);
  \draw[flowteal] (disease.east)--(router.west);
  \draw[flowwarm] (target.east) to[out=0,in=200] (router.south west);
  \node[captionnote,text width=3.40cm] at (5.03,-4.36)
    {Patient, disease and target context};

  % The three rows are categories, not a claim that only three sources exist.
  \node[stageblue,minimum width=3.58cm] (shared) at (9.45,-1.97)
    {$z_i$ \quad shared path};
  \node[stageteal,minimum width=3.58cm] (dsrc) at (9.45,-3.08)
    {$\sum_d\pi_{i,k,d}\,a_d(z_i)$ \quad disease};
  \node[stagewarm,minimum width=3.58cm] (osrc) at (9.45,-4.19)
    {$\sum_o\pi_{i,k,o}\,a_o(z_i)$ \quad outcome};
  \draw[flowblue] (router.east) to[out=25,in=180] (shared.west);
  \draw[flowteal] (router.east)--(dsrc.west);
  \draw[flowwarm] (router.east) to[out=-25,in=180] (osrc.west);

  \node[stage,minimum width=2.33cm] (mix) at (14.07,-2.56)
    {weighted residual\\$z_{i,k}$};
  \node[outcome,minimum width=2.33cm] (hazards) at (14.07,-4.02)
    {interval hazards\\$\lambda_{i,k,t}$};
  \draw[flowblue] (shared.east) to[out=0,in=155] (mix.north west);
  \draw[flowteal] (dsrc.east)--(mix.west);
  \draw[flowwarm] (osrc.east) to[out=0,in=205] (mix.south west);
  \draw[flowwarm] (mix.south)--(hazards.north);

  \draw[divider] (0,-5.34)--(15.84,-5.34);
  \node[captionnote,anchor=west,text width=15.40cm,align=left]
    at (.20,-5.68)
    {The shared path remains unscaled; route weights govern adapter residuals. The bank includes eligible borrowed sources, whose adapter parameters are detached.};
\end{tikzpicture}
}
  \caption{\textbf{Source-anchored disease--outcome routing.} For each patient
  and target outcome, the router mixes a shared representation with residual
  adapters indexed by fine disease and outcome. Its weights depend on the
  patient, disease and target. Sources with no training support are masked and
  low-event sources are shrunk. An adapter is updated by its own eligible
  supervision, while borrowed adapter parameters are detached. The mixed
  representation enters a shared linear head for interval hazards and
  cumulative incidence. This is a model schematic, not an empirical route-weight
  or performance plot.}
  \label{fig:source-router}
\end{figure}

\begin{figure}[p]
  \resultfigure[\textwidth]{supp_routing}{Routing}{Routed versus joint
    and adapter-bank comparison.}
  \caption{\textbf{Routed versus joint model, and adapter-bank comparisons.}
  \textbf{A}, Paired
  one-year AUC gain of the routed over the joint model by clinical domain, with the prespecified margin of practical equivalence shaded. \textbf{B},
  Panel gains of adapter banks trained under the same protocol: no adapters
  versus the joint model (a check that the architecture alone changes nothing),
  disease adapters, outcome adapters, each given the other, both, and their
  interaction. The router configuration was chosen on the tuning set.}
  \label{fig:routing}
\end{figure}

\begin{figure}[p]
  \resultfigure[0.9\textwidth]{supp_rescue}{Routing where joint learning
    reduced performance}{Cells harmed on tuning and their test-period gains.}
  \caption{\textbf{Routing in cells where joint learning reduced performance
  on the tuning set.} Cells were classified on the tuning set by whether the joint model
  was worse than the specialist. \textbf{A}, Routed-versus-joint test gain
  against the tuning-set harm. \textbf{B}, Test-period gains of the joint and
  routed models over the specialist in each stratum.}
  \label{fig:supp-rescue}
\end{figure}

\begin{figure}[p]
  \resultfigure[0.9\textwidth]{supp_router_usage}{Router weights}{Router
    weights by disease and clinical domain.}
  \caption{\textbf{Router weights by disease and clinical domain.} \textbf{A}, Mean
  routing weight on each disease adapter by patients' disease (the patient's
  own disease hatched). \textbf{B}, Outcome-adapter weights by clinical domain.
  \textbf{C}, Share of the routed correction carried by the disease and
  outcome banks. \textbf{D}, Routing intensity across patients.}
  \label{fig:supp-router-usage}
\end{figure}

\begin{figure}[p]
  \resultfigure[\textwidth]{supp_strategy}{Joint model versus per-cell
    selection}{Joint model against the specialist and the tuning-selected best model.}
  \caption{\textbf{Joint model versus the best model selected for each cell.}
  \textbf{A}, Test AUC of the joint model against the specialist in each cell, with point size reflecting training events. \textbf{B}, Against the best of
  the specialist, disease-pooled and outcome-pooled models chosen per cell on
  the tuning set. \textbf{C}, Which model was chosen, by clinical domain.}
  \label{fig:strategy}
\end{figure}

\begin{figure}[p]
  \resultfigure[0.6\textwidth]{supp_national_where}{Effect of national
    pretraining by subgroup}{Effects by vocabulary coverage, support and year.}
  \caption{\textbf{Effect of national pretraining by vocabulary coverage,
  training support and year of treatment.} Patient-weighted
  improvement in one-year Brier score ($\times$100, 95\% CI) of the joint model
  initialised from the national encoder, used directly or after local
  adaptation, over the locally pretrained joint model, by the share of a
  patient's history covered by the national vocabulary, by the number of
  training events of the predicted combination, and by year of treatment.}
  \label{fig:supp-national-where}
\end{figure}

\begin{figure}[p]
  \resultfigure[0.6\textwidth]{fig_who_benefits}{Gain by patient subgroup}
    {Paired Brier improvement of the joint model over the specialist and over
    XGBoost by patient subgroup.}
  \caption{\textbf{Gain from joint learning by patient subgroup.} Patient-weighted
  improvement in one-year Brier score ($\times$100, 95\% bootstrap CI) of the
  joint model over the specialist and over XGBoost, pooled across cells, by
  subgroups defined at the start of treatment. Positive values favour the
  joint model.}
  \label{fig:who-benefits}
\end{figure}

\begin{figure}[p]
  \resultfigure[0.9\textwidth]{supp_calibration_maps}{Calibration maps}{Local
    calibration over the patient representation.}
  \caption{\textbf{Local calibration across the patient representation.} Two-dimensional projection of test patients' representations
  for example outcomes: event status, routing displacement, kernel-smoothed
  predicted risk and the IPCW-weighted local residual (observed minus
  predicted).}
  \label{fig:supp-calibration-maps}
\end{figure}

\begin{figure}[p]
  \resultfigure{fig_patient_map}{Patient representation}
    {A: disease groups on the joint model's patient map. B--F: observed risk of
    five outcomes on the same map. G: local calibration.}
  \caption{\textbf{Observed outcome risk across the joint model's patient
  representation.} \textbf{A}, Test
  patients on a two-dimensional projection of the joint model's patient
  representation (fitted on tuning patients), with the densest half of each
  disease group outlined. \textbf{B--F}, The same map coloured by the observed
  one-year risk of five outcomes from different clinical domains, estimated
  with censoring-weighted kernels and shown relative to each outcome's average. Regions with too few evaluable patients are blank. \textbf{G}, Observed minus
  predicted risk of the joint model for the outcome in \textbf{B}. All 69
  outcomes are shown in \cref{fig:supp-embedding-gallery-1}--\cref{fig:supp-embedding-gallery-6}.
  The map is descriptive and carries no performance claim.}
  \label{fig:patient-map}
\end{figure}

\begin{figure}[p]
  \resultfigure{supp_embedding_comparison}{Patient representations across training stages}
    {The same test patients under five model stages, from local and national
    pretraining to the fine-tuned joint models.}
  \caption{\textbf{Patient representations at each training stage.}
  The same test patients under five model stages: pretraining on the cohort,
  the joint model from that encoder, national pretraining, national pretraining
  adapted to the cohort, and the nationally initialised joint model. The first
  row outlines where each disease group lies, and the others show the local mean
  (Gaussian kernel) of age at treatment, the number of recorded events before
  treatment (log scale) and the IPI category among patients with a recorded IPI,
  left blank where too few patients are nearby. Each column's projection
  was fitted on tuning patients and applied unchanged to test patients, so
  coordinates are comparable within a column but not between columns. The maps
  are descriptive and carry no performance claim.}
  \label{fig:supp-embedding-comparison}
\end{figure}

\begin{figure}[p]
  \resultfigure[0.95\textwidth]{supp_embedding_gallery_1}{Outcome risk across the representation (1/6)}
    {Observed one-year risk of each outcome on the joint model's patient map.}
  \caption{\textbf{Observed one-year risk of each outcome across the patient representation (1 of 6).}
  Every panel shows the same two-dimensional projection of the joint model's
  patient embedding (fitted on tuning patients) for test patients. Colour is
  the censoring-weighted kernel estimate of the observed one-year risk of one
  outcome, relative to that outcome's average (shown under its name): bright
  regions carry up to four times the average risk and dark regions a quarter.
  Regions with too few evaluable patients are blank. Outcomes are grouped by
  clinical domain, named above each group.}
  \label{fig:supp-embedding-gallery-1}
\end{figure}

\begin{figure}[p]
  \resultfigure[0.95\textwidth]{supp_embedding_gallery_2}{Outcome risk across the representation (2/6)}
    {Observed one-year risk of each outcome on the joint model's patient map.}
  \caption{\textbf{Observed one-year risk of each outcome across the patient representation (2 of 6).}
  Every panel shows the same two-dimensional projection of the joint model's
  patient embedding (fitted on tuning patients) for test patients. Colour is
  the censoring-weighted kernel estimate of the observed one-year risk of one
  outcome, relative to that outcome's average (shown under its name): bright
  regions carry up to four times the average risk and dark regions a quarter.
  Regions with too few evaluable patients are blank. Outcomes are grouped by
  clinical domain, named above each group.}
  \label{fig:supp-embedding-gallery-2}
\end{figure}

\begin{figure}[p]
  \resultfigure[0.95\textwidth]{supp_embedding_gallery_3}{Outcome risk across the representation (3/6)}
    {Observed one-year risk of each outcome on the joint model's patient map.}
  \caption{\textbf{Observed one-year risk of each outcome across the patient representation (3 of 6).}
  Every panel shows the same two-dimensional projection of the joint model's
  patient embedding (fitted on tuning patients) for test patients. Colour is
  the censoring-weighted kernel estimate of the observed one-year risk of one
  outcome, relative to that outcome's average (shown under its name): bright
  regions carry up to four times the average risk and dark regions a quarter.
  Regions with too few evaluable patients are blank. Outcomes are grouped by
  clinical domain, named above each group.}
  \label{fig:supp-embedding-gallery-3}
\end{figure}

\begin{figure}[p]
  \resultfigure[0.95\textwidth]{supp_embedding_gallery_4}{Outcome risk across the representation (4/6)}
    {Observed one-year risk of each outcome on the joint model's patient map.}
  \caption{\textbf{Observed one-year risk of each outcome across the patient representation (4 of 6).}
  Every panel shows the same two-dimensional projection of the joint model's
  patient embedding (fitted on tuning patients) for test patients. Colour is
  the censoring-weighted kernel estimate of the observed one-year risk of one
  outcome, relative to that outcome's average (shown under its name): bright
  regions carry up to four times the average risk and dark regions a quarter.
  Regions with too few evaluable patients are blank. Outcomes are grouped by
  clinical domain, named above each group.}
  \label{fig:supp-embedding-gallery-4}
\end{figure}

\begin{figure}[p]
  \resultfigure[0.95\textwidth]{supp_embedding_gallery_5}{Outcome risk across the representation (5/6)}
    {Observed one-year risk of each outcome on the joint model's patient map.}
  \caption{\textbf{Observed one-year risk of each outcome across the patient representation (5 of 6).}
  Every panel shows the same two-dimensional projection of the joint model's
  patient embedding (fitted on tuning patients) for test patients. Colour is
  the censoring-weighted kernel estimate of the observed one-year risk of one
  outcome, relative to that outcome's average (shown under its name): bright
  regions carry up to four times the average risk and dark regions a quarter.
  Regions with too few evaluable patients are blank. Outcomes are grouped by
  clinical domain, named above each group.}
  \label{fig:supp-embedding-gallery-5}
\end{figure}

\begin{figure}[p]
  \resultfigure[0.95\textwidth]{supp_embedding_gallery_6}{Outcome risk across the representation (6/6)}
    {Observed one-year risk of each outcome on the joint model's patient map.}
  \caption{\textbf{Observed one-year risk of each outcome across the patient representation (6 of 6).}
  Every panel shows the same two-dimensional projection of the joint model's
  patient embedding (fitted on tuning patients) for test patients. Colour is
  the censoring-weighted kernel estimate of the observed one-year risk of one
  outcome, relative to that outcome's average (shown under its name): bright
  regions carry up to four times the average risk and dark regions a quarter.
  Regions with too few evaluable patients are blank. Outcomes are grouped by
  clinical domain, named above each group.}
  \label{fig:supp-embedding-gallery-6}
\end{figure}

\begin{figure}[p]
  \resultfigure{supp_attribution_atlas}{What the model relies on}{Differences between patients by source and time, and findings across every outcome.}
  \caption{\textbf{What the model relies on, for every outcome.} Rows are
  outcomes grouped by clinical domain. Integrated gradients split each held-out
  patient's predicted log-odds at one year into a reference, which keeps the
  timing of care, age and subtype but replaces every recorded event by an
  average event, and the contribution of each event. \textbf{A}, The variance of
  the predictions between patients, divided into the part carried by the timing
  of care alone (grey) and the parts carried by what was recorded, by source.
  Shares are signed and add up to 100\% per outcome. \textbf{B}, The same
  division by time before treatment. \textbf{C}, Mean change in predicted
  log-odds among the patients who have each finding, shown as the corresponding
  odds ratio. Boxed cells are the
  prespecified expected effects, dots mark 95\% patient-bootstrap intervals that
  exclude zero, and grey cells had fewer than 20 patients with the finding.}
  \label{fig:attribution-atlas}
\end{figure}

\begin{figure}[p]
  \resultfigure{supp_attribution_structure}{Outcome-specific use of the history}{Agreement between outcomes and what each part of the record adds.}
  \caption{\textbf{The model uses the same history differently for different
  outcomes.} \textbf{A}, For every pair of outcomes, how much they agree on which
  episodes of a patient's history raise or lower risk (Spearman correlation of
  episode contributions within patients, averaged over patients). Episodes are
  events within one day of each other, and each clinical domain is outlined.
  \textbf{B}, Change in held-out AUC ($\times$100) when the same fitted model
  reads only part of the record, averaged over all outcomes and over the
  outcomes of each domain. The model is not retrained, so this measures what it
  relies on, not what a model trained without that input could achieve.}
  \label{fig:attribution-structure}
\end{figure}

\begin{figure}[p]
  \resultfigure{supp_attribution_patient}{One patient's history, read six ways}{For a synthetic patient, the episodes that move each of six risks, at their time.}
  \caption{\textbf{One synthetic patient's history read for six
  complications.} Each track is one outcome from a different clinical domain.
  Circles are the episodes of care (events within one day of each other) that
  move its predicted risk, at their time before treatment, sized by the change
  in predicted log-odds and coloured by direction (sienna raises, indigo
  lowers). The episodes that matter most for any outcome are labelled with
  their main findings and timing. The predicted one-year risk is given at the
  right, with the risk for the same timing of care and average findings. For
  every patient, the episodes that most raise and lower each risk can be
  extracted in the same way and shown beside the prediction. The patient and
  its contributions are invented, and records of real patients stay in the
  secure environment.}
  \label{fig:attribution-patient}
\end{figure}
